\section{Experiments on real stacks}\label{sec:experiments}

All experiments run on one NVIDIA H100 NVL (94\,GB). \Cref{app:details} gives the full configurations. \todo{link an anonymized repository with code and raw outputs.}

\subsection{UAQ leaves quantized rollouts unchanged}\label{sec:s1}

\paragraph{Protocol.} We test \cref{thm:invariance} with the kernels QuRL would use on Hopper GPUs, served by vLLM 0.33 \citep{kwon2023vllm}.
\begin{itemize}[leftmargin=*,itemsep=1pt]
\item \emph{Models:} Qwen2.5-0.5B-Instruct \citep{qwen2024qwen25}, and Qwen2.5-Math-1.5B and -7B \citep{yang2024qwen25math}.
\item \emph{Checkpoints:} the original and a UAQ checkpoint, both stored in BF16 like the original.
\item \emph{Rollout quantizers:} INT8 W8A8 with vLLM's CUTLASS INT8 GEMM, FP8 per-tensor, and FP8 with per-channel weights and per-token activations (both on CUTLASS FP8 GEMMs), plus unquantized BF16.
\item \emph{Tokens:} 512 GSM8K \citep{cobbe2021gsm8k} responses sampled from the BF16 model at temperature 1. Every token (178{,}668, 193{,}103 and 216{,}979 for the three sizes) is teacher-forced under each configuration.
\item \emph{Noise floor:} the same configuration re-run with shuffled batches.
\item \emph{Reference policy:} because UAQ re-rounds $W/s$ and $s\gamma$ in BF16, each quantized checkpoint is compared with the BF16 policy of the \emph{same} checkpoint.
\item \emph{Metrics:} the primary metric is the mean $|\log\mu-\log\pi|$ over sampled tokens; we also report the $k_3$ estimate of $\KL(\pi\|\mu)$. Confidence intervals come from a paired bootstrap over responses.
\end{itemize}

\begin{table}[t]
\centering\footnotesize
\begin{tabularx}{\linewidth}{@{}L l l l l l@{}}
\toprule
& & \multicolumn{2}{c}{Identity, mean $|\Delta\log p|$; identical tokens} & \multicolumn{2}{c}{$s=1.5$: ratio UAQ / original, 95\% CI}\\
\cmidrule(lr){3-4}\cmidrule(l){5-6}
Rollout quantizer & Model & shuffled re-run & UAQ $s=2$ vs.\ original & mean $|\Delta\log p|$ & KL ($k_3$)\\
\midrule
\textbf{INT8 W8A8} (QuRL) & 0.5B & $1.4\times10^{-4}$; 99.69\% & \textbf{0}; 100\% & 0.994 [0.984, 1.004] & 0.978 [0.951, 1.007]\\
 & 1.5B & $7.7\times10^{-5}$; 99.75\% & \textbf{0}; 100\% & 0.996 [0.988, 1.005] & 0.994 [0.966, 1.022]\\
 & 7B & 0; 100\% & \textbf{0}; 100\% & 1.013 [1.004, 1.021] & 1.022 [0.995, 1.049]\\
\midrule
FP8 per-tensor & 0.5B & $1.2\times10^{-4}$; 99.78\% & \textbf{0}; 100\% & 1.004 [0.995, 1.014] & 1.013 [0.985, 1.043]\\
 & 1.5B & $9.7\times10^{-5}$; 99.70\% & $2.7\times10^{-5}$; 99.95\%$^\dagger$ & 0.976 [0.968, 0.984] & 0.923 [0.888, 0.956]\\
 & 7B & 0; 100\% & \textbf{0}; 100\% & 0.946 [0.939, 0.954] & 0.862 [0.830, 0.891]\\
\midrule
FP8 per-channel/token & 0.5B & $1.1\times10^{-4}$; 99.78\% & \textbf{0}; 100\% & 1.007 [0.999, 1.016] & 1.013 [0.989, 1.039]\\
 & 1.5B & $6.1\times10^{-5}$; 99.75\% & $1.7\times10^{-4}$; 99.77\%$^\dagger$ & 0.968 [0.960, 0.976] & 0.888 [0.854, 0.921]\\
 & 7B & 0; 100\% & \textbf{0}; 100\% & 0.943 [0.935, 0.951] & 0.964 [0.936, 0.997]\\
\bottomrule
\end{tabularx}
\caption{UAQ on vLLM's quantized kernels. An $s$-fold error reduction, as claimed, would give ratios of 0.667 for mean $|\Delta\log p|$ and 0.444 for KL. $^\dagger$Run-to-run variation of vLLM's FP8 path at 1.5B: in same-order repeats an $s=2$ run reproduced an original run bit for bit.}\label{tab:s1}
\end{table}

\paragraph{Results} (\cref{tab:s1}).
\begin{itemize}[leftmargin=*,itemsep=2pt]
\item \textbf{At $s=2$ the invariance is exact on the hardware kernel.} QuRL's INT8 W8A8 rollout gives bit-identical log-probabilities at all three sizes, as does FP8 at 0.5B and 7B. No INT8 or FP8 code changes, and the stored INT8 scales are exactly halved.
\item \textbf{At $s=1.5$, UAQ does not reduce QuRL's INT8 mismatch at any size.}
  \begin{itemize}[leftmargin=*,itemsep=1pt]
  \item The ratios are 0.994--1.013 for mean $|\Delta\log p|$ and 0.978--1.022 for KL, against 0.667 and 0.444 for the claimed reduction.
  \item This holds although BF16 storage of $W/1.5$ re-rounds 4.0--4.3\% of the INT8 codes.
  \end{itemize}
\item \textbf{FP8 shows small changes at 1.5B and 7B (ratios 0.94--0.98), never near $s$-fold.}
  \begin{itemize}[leftmargin=*,itemsep=1pt]
  \item \Cref{thm:invariance} predicts no systematic change. Both FP8 modes of a model share one BF16 re-rounding of $W$ and $\gamma$, so the table holds three draws, two of them favourable.
  \item The per-token errors before and after UAQ correlate only 0.43--0.64, and $|\log\mu-\log\pi|$ shrinks on about as many tokens as it grows.
  \end{itemize}
\end{itemize}

\paragraph{By-products that matter for rollout research} (not specific to UAQ).
\begin{itemize}[leftmargin=*,itemsep=1pt]
\item vLLM's default engine is not batch-invariant at 0.5B and 1.5B: shuffled batching changes 0.22--0.31\% of tokens, BF16 included.
\item The mismatch is heavy-tailed: the top 1\% of tokens carry 38--48\% of the KL.
\item The INT8-to-FP8 KL ratio grows with model size, from 1.0--1.1$\times$ at 0.5B to 2.8--3.0$\times$ at 7B. This is consistent with emergent activation outliers \citep{dettmers2022llmint8}.
\end{itemize}

\subsection{Matched optimizer settings reproduce UAQ exactly}\label{sec:s2}

\paragraph{Protocol.} We use $s=2$, where every rescaling is exact, so \cref{cor:pipeline} predicts bit-identical training. We run a QuRL-style loop in PyTorch on Qwen2.5-0.5B-Instruct and GSM8K.
\begin{itemize}[leftmargin=*,itemsep=1pt]
\item \emph{Loop:} GRPO advantages \citep{shao2024deepseekmath}, the decoupled PPO loss with TIS ($C=2$) and ACR, and AdamW ($\lambda=0.01$, $\varepsilon=10^{-8}$) with global-norm clipping at $\tau=0.1$.
\item \emph{Rollouts:} INT8 W8A8 rollouts, with per-channel weights and per-token activations, sampled from seeded uniforms.
\item \emph{Training:} BF16 autocast with FP32 master weights, under deterministic kernels.
\item \emph{Length:} 20 rollout batches of $8\times4$ responses with 4 updates each, 80 in all, at $\eta=10^{-6}$ and $10^{-5}$.
\item \emph{Arms:} UAQ with plain AdamW (R4); an R4 repeat; the matched settings of \cref{thm:reparam}(b)--(c) without UAQ (R5); R5 with ordinary clipping (R5o); learning rates matched but not $\varepsilon$, weight decay or clipping (R5x); and plain QuRL (R3).
\item \emph{Comparison:} each step is fingerprinted (SHA-256) on its sampled tokens, $\mu$, $\pi_{\mathrm{old}}$, and the weights mapped to original coordinates.
\end{itemize}

\begin{table}[t]
\centering\footnotesize
\begin{tabularx}{\linewidth}{@{}l L l L l@{}}
\toprule
Arm & Differs from R4 in & Steps identical to R4 & First difference & Max.\ weight diff.\ after step 20\\
\midrule
R4 repeat & nothing & 20/20; 20/20 & none & 0; 0\\
\textbf{R5} & coordinates only (\cref{thm:reparam}(b)--(c)) & \textbf{20/20; 20/20} & \textbf{none} & \textbf{0; 0}\\
\midrule
R5o & clipping norm in original coordinates & 0/20; 0/20 & weights after step 1 & $1.63\times10^{-4}$; $1.13\times10^{-3}$\\
R5x & learning rates only & 0/20; 0/20 & weights after step 1 & $1.65\times10^{-4}$; $1.32\times10^{-3}$\\
R3 & no UAQ, plain settings & 0/20; 0/20 & weights after step 1 & $1.36\times10^{-4}$; $1.12\times10^{-3}$\\
\bottomrule
\end{tabularx}
\caption{Exact test of \cref{cor:pipeline} at $s=2$ ($\eta=10^{-6}$; $10^{-5}$). A step counts as identical when all four SHA-256 fingerprints equal R4's. Weight differences are in original coordinates.}\label{tab:s2}
\end{table}

\paragraph{Results} (\cref{tab:s2}).
\begin{itemize}[leftmargin=*,itemsep=2pt]
\item \textbf{R5 reproduces UAQ bit for bit.} It matched R4 in every fingerprint at every step, at both learning rates.
  \begin{itemize}[leftmargin=*,itemsep=1pt]
  \item Its weights after steps 1 and 20 equal R4's mapped weights in every element, and so do the gradient norms that set the clip.
  \item Clipping bound on every update with a nonzero gradient: 70 of 70 at $\eta=10^{-6}$ and 44 of 44 at $10^{-5}$.
  \item Weight decay was active at $10^{-5}$. At $10^{-6}$ it changes no FP32 weight at all, since $\eta\lambda<2^{-25}$.
  \end{itemize}
\item \textbf{Rule (c) is necessary whenever clipping binds.}
  \begin{itemize}[leftmargin=*,itemsep=1pt]
  \item On the first update, the gradient norm is 1.131 in UAQ coordinates and 0.874 in original coordinates. Ordinary clipping (R5o) therefore leaves the gradient $1.29\times$ larger than UAQ's clip does.
  \item R5o diverged after the first update, as early as plain QuRL (R3). So did matching the learning rates alone (R5x).
  \end{itemize}
\item \textbf{The equivalence also holds in VeRL's training engine.} We implemented the matched settings and rule (c) for VeRL's FSDP2 actor (\cref{app:verl}) and ran the same comparison inside that engine on the H100: $s=2$, three updates on one GSM8K batch, $\eta=10^{-5}$, clipping at 0.1, BF16 mixed precision, strict deterministic algorithms.
  \begin{itemize}[leftmargin=*,itemsep=1pt]
  \item R5 and an R4 repeat were bit-identical to R4 in all 494M parameters, with identical gradient norms (3.859, 2.573, 3.110).
  \item R5 with ordinary clipping differed in 99.98\% of the Q/K/V weights; its first gradient norm was 2.791 against 3.859.
  \end{itemize}
\end{itemize}

\subsection{Update capture on a real PPO step}\label{sec:capture}

\Cref{prop:capture} predicts that a sub-grid update flips a code with probability $|d|/\Delta$. We checked this in VeRL with INT8 W8A8 rollouts that re-quantize the FP32 master weights at every weight sync (\cref{app:verl}). Before training, the codes and scales are bit-identical to an independent quantization of the checkpoint.

After the first PPO step on Qwen2.5-0.5B-Instruct (four AdamW updates at $\eta=10^{-6}$), 0.21\% of layer~0's Q/K/V codes had flipped.
\begin{itemize}[leftmargin=*,itemsep=1pt]
\item \emph{Prediction:} we estimate $|d|$ for each row from the move of its absmax element (median $2\times10^{-6}$). \Cref{prop:capture} then predicts 0.20\%, or 0.24\% when the shift of the grid itself is counted.
\item \emph{By projection,} observed vs.\ predicted: Q 0.16\% vs.\ 0.15\%, K 0.29\% vs.\ 0.18\%, V 0.53\% vs.\ 0.58\%.
\item \emph{Row by row} the agreement is loose (correlation 0.2--0.4), because $|d|$ is estimated from a single element per row.
\end{itemize}
This is an aggregate check on one layer and one step. It confirms the mechanism: re-quantization captures sub-grid updates at the predicted rate, so the rollout model is not frozen.

\subsection{GSM8K PPO: a reproduction attempt and the UAQ experiment}\label{sec:s3}

\paragraph{Setup.} We use QuRL's GSM8K PPO configuration in VeRL v0.9.1 \citep{sheng2025hybridflow}, with vLLM 0.24 rollouts and Qwen2.5-0.5B-Instruct.
\begin{itemize}[leftmargin=*,itemsep=1pt]
\item \emph{PPO:} a critic with GAE ($\gamma=\lambda=1$), batch 256, responses up to 512 tokens, and 15 epochs, i.e.\ 435 steps.
\item \emph{Regularizers:} KL penalty 0.001 in the reward and entropy bonus 0.001, as in VeRL's documented PPO baseline for this model (56.7\% at actor learning rate $10^{-6}$). QuRL states neither.
\item \emph{Validation:} greedy, on the 1{,}319 test problems, every 29 steps.
\item \emph{QuRL arms:} INT8 W8A8 rollouts re-quantized at every sync, token-level TIS with $C=2$, and ACR.
\end{itemize}

\paragraph{QuRL's Table~1 does not reproduce at the stated learning rate.} QuRL reports 55.35\% for BF16 RL and 53.55\% for QuRL INT8 at ``the relatively high learning rate ($10^{-5}$)''. At actor learning rate $10^{-5}$ (critic $10^{-5}$), our BF16 and INT8 runs reached 0.0\% and 2.9\% after 435 steps; their best of 15 validations was 0.0\% and 3.4\%.
\begin{itemize}[leftmargin=*,itemsep=2pt]
\item \textbf{Neither run learned the task.}
  \begin{itemize}[leftmargin=*,itemsep=1pt]
  \item The BF16 run answered only one more prompt correctly after step 4, and its responses shrank to 2--3 tokens.
  \item The INT8 run settled on answers of about 9 tokens, and later on near-deterministic outputs that ran to the 512-token cap.
  \end{itemize}
\item \textbf{In our configuration, the learning rate alone separates learning from collapse.} Ten-step pilots in the same configuration at $10^{-6}$ isolate it. After 10 steps:
  \begin{itemize}[leftmargin=*,itemsep=1pt]
  \item the KL from the reference model was $2.1\times10^{-3}$ per token at $10^{-6}$ against 0.15 at $10^{-5}$;
  \item the PPO clip fraction was 0.001 against 0.010--0.022;
  \item the entropy stayed at 0.44--0.53 at $10^{-6}$, but rose from 0.50 to 0.80 at $10^{-5}$.
  \end{itemize}
\item \textbf{Mechanism.} At step 1, the critic's newly initialized value head predicts values from $-20$ to $+33$ for returns in $[0,1]$, so the early advantages mostly reflect its error.
  \begin{itemize}[leftmargin=*,itemsep=1pt]
  \item At $10^{-5}$, the actor follows that error ten times as fast and loses the rare correctly formatted answers before the critic becomes useful.
  \item Once no answer is rewarded, the per-token KL penalty is the only reward left, and it shortens the responses.
  \end{itemize}
\item \textbf{What the paper's run must have done differently.} QuRL does not state the critic's learning rate, warmup, KL penalty or entropy bonus. Two explanations fit:
  \begin{enumerate}[leftmargin=*,itemsep=1pt]
  \item the actor ran at $10^{-6}$ and $10^{-5}$ is the critic's rate, as in VeRL's baseline, from whose 56.7\% QuRL's 55.35\% differs by 1.4 points;
  \item the run used a stabilizer the paper does not state.
  \end{enumerate}
  At $10^{-6}$ our configuration trains: QuRL INT8 with seed 1 reaches 41.4\% at step 29 and 50.1\% at step 435.
\end{itemize}

\paragraph{The INT8 mismatch grows although the weight error cannot.} In that run, the sampled-token mismatch ($k_3$) rose from $1.1\times10^{-2}$ at step 1 to window medians of $7.7$--$8.1\times10^{-2}$ over steps 60--175. It then fell to $3.6\times10^{-2}$, still about 65$\times$ the BF16 floor.
\begin{itemize}[leftmargin=*,itemsep=1pt]
\item It tracked neither the policy's entropy (Spearman 0.13) nor its KL from the reference model ($-0.22$).
\item By \cref{prop:persistence}, the weight rounding error is stationary. The growth must therefore come from sensitivity, activations or contexts (\cref{sec:tracking}).
\item Attributing it requires per-layer error and activation logging, which we leave to future work.
\end{itemize}

\paragraph{The UAQ experiment (pre-registered).} We compare six arms at $\eta=10^{-6}$ and $s=1.5$.
\begin{itemize}[leftmargin=*,itemsep=1pt]
\item \emph{Arms:}
  \begin{itemize}[leftmargin=*,itemsep=1pt]
  \item R1, BF16 rollouts; R2, BF16 rollouts with UAQ;
  \item R3, QuRL INT8; R4, QuRL with UAQ;
  \item R5, QuRL with the settings of \cref{thm:reparam}(b)--(c);
  \item R6, QuRL with the global learning rate $\times1.5$, QuRL's own control.
  \end{itemize}
\item \emph{Implementation:} all arms share one FSDP2 code path and differ only in rollout precision, starting checkpoint and optimizer settings.
\item \emph{Seeds:} 1--3 for R1, R2, R3 and R6, and 1--5 for R4 and R5. The arms of one seed share prompt order and sampling seed.
\end{itemize}
The hypotheses are:
\begin{itemize}[leftmargin=*,itemsep=1pt]
\item \textbf{H1:} R4 and R5 are equal in distribution (\cref{cor:pipeline});
\item \textbf{H2:} UAQ's effect is not specific to quantization, $(\text{R2}-\text{R1})\approx(\text{R4}-\text{R3})$;
\item \textbf{H3:} QuRL's control does not match UAQ, $\text{R6}\neq\text{R4}$.
\end{itemize}
The primary endpoint is the mean test accuracy over the last five validations. H1 and H2 are equivalence claims, tested by two one-sided tests with a margin of $\pm2.0$ points. H3 is tested for a difference with a 95\% interval. Analyses, margins and decision rules were fixed and their scripts hashed before the relevant arms were seen (\cref{app:prereg}).

\todo{Pre-registered results: verdicts for H1--H3, per-arm late accuracy with intervals, and training curves (figure).}

\paragraph{\pending{Preliminary observation (seed 1 only; to be replaced by the pre-registered analysis).}}
Training at $10^{-6}$ is not immune to the early drift seen at $10^{-5}$.
\begin{itemize}[leftmargin=*,itemsep=1pt]
\item In R5 with seed 1, the few correct answers vanished within three steps, and no answer was rewarded until step 374.
\item The KL penalty then pulled the policy back toward the reference model, and the run learned, reaching 46.6\% at step 435.
\item R3 and R4 with the same seed had a correct answer in every batch.
\end{itemize}
By \cref{cor:pipeline}, R4 and R5 run the same algorithm. A single stall is therefore an outcome of the configuration, not evidence against H1. Stall counts per arm are part of the pre-registered exploratory analysis.
